\documentclass{beamer}
%
\mode<presentation>
{
  \usetheme{default}
  \usecolortheme{default}
  \usefonttheme{default}
  \setbeamertemplate{navigation symbols}{}
  \setbeamertemplate{caption}[numbered]
}

\usepackage[english]{babel}
\usepackage[utf8]{inputenc}
\usepackage[T1]{fontenc}
\usepackage{lmodern}
\usepackage{amsmath}
\usepackage{amssymb}
\usepackage{graphicx}

\title[Reference Presentation 2: Finite Blocklength MIMO Precoding]
{Finite Blocklength MIMO Precoding\\
With Mixed Power and QoS Constraints}

\author{Shanik Hubert}

\institute{
Paper Explanation Presentation\\[0.2cm]
\small Based on: Jingchen Wu et al., \textit{IEEE Communications Letters}, 2026
}

\date{}



\begin{document}

\begin{frame}{}
    \titlepage
\end{frame}

\begin{frame}{Contents}
  \tableofcontents
\end{frame}

\section{Overall Idea of the Paper}
%------------------------------------------------
\begin{frame}{Layer 0: Overall Idea of the Paper}
\small

\begin{block}{Core Problem}
Design downlink MIMO precoders that maximize finite-blocklength weighted sum rate while satisfying mixed power constraints and QoS requirements.
\end{block}

\begin{itemize}
    \item To reduce difficulty of precoder design problem, it is reformulated into an easier SINR allocation problem.\
    
    \[
        \text{Hard precoder design}
        \quad \Longrightarrow \quad
        \text{SINR allocation problem}
        \]
\end{itemize}
        
\begin{itemize}
    \item Instead of directly optimizing the precoding matrix \(V\), the paper optimizes target SINRs \(\gamma_k\).
    \item A power feasibility function \(P(\gamma)\) checks whether the target SINRs can be achieved.
    \item Two algorithms are proposed:
    \begin{itemize}
        \item BDGM: bi-direction gradient method.
        \item IWF: iterative water-filling algorithm.
    \end{itemize}
\end{itemize}

\end{frame}
%------------------------------------------------

%------------------------------------------------
\section{Layer 1: Main Research Gap}
\begin{frame}{Layer 1: Main Research Gap}
\small

\begin{itemize}
    \item Existing finite-blocklength precoding works usually focus on limited settings.

    \item Many methods consider only a single sum power constraint:
    \[
        \sum_{k=1}^{K}\|v_k\|^2 \le P
    \]

    \item Practical MIMO systems may have different power limits:
    \begin{itemize}
        \item sum power constraint,
        \item per-antenna power constraint,
        \item per-group or per-AP power constraint.
    \end{itemize}

    \item Therefore, a general finite-blocklength precoding framework is needed for mixed power constraints and QoS guarantees.
\end{itemize}

\end{frame}

\section{Layer 2: Downlink MIMO System Model}
%------------------------------------------------
%------------------------------------------------
\begin{frame}{Layer 2: Downlink MIMO System Model}
\small

\begin{itemize}
    \item The system has \(N\) transmit antennas serving \(K\) single-antenna users.

    \item The transmitted signal is:
    \[
        \mathbf{x}
        =
        \sum_{k=1}^{K}
        \mathbf{D}_k\mathbf{v}_k s_k
    \]

    \item Variables:
    \begin{itemize}
        \item \(\mathbf{x}\in\mathbb{C}^{N}\): transmitted signal vector from all antennas.
        \item \(N\): number of transmit antennas.
        \item \(K\): number of users.
        \item \(k\): user index, where \(k=1,\ldots,K\).
        \item \(s_k\): data symbol intended for user \(k\).
        \item \(\mathbb{E}[s_k]=0\), \(\mathbb{E}[|s_k|^2]=1\): symbol has zero mean and unit power.
        \item \(\mathbf{v}_k\in\mathbb{C}^{N}\): precoding vector for user \(k\).
    \end{itemize}

    \item Antenna-user association:
    \[
        \mathbf{D}
        =
        [\mathbf{d}_1,\ldots,\mathbf{d}_K]
        \in \{0,1\}^{N\times K}
    \]

    \[
        [\mathbf{D}]_{nk}
        =
        \begin{cases}
        1, & \text{if antenna } n \text{ serves user } k,\\
        0, & \text{otherwise}.
        \end{cases}
    \]

    \item \(\mathbf{D}_k=\operatorname{diag}(\mathbf{d}_k)\in\{0,1\}^{N\times N}\) selects only the antennas assigned to user \(k\).
\end{itemize}

\end{frame}
%------------------------------------------------
\begin{frame}{Layer 2: Received Signal Model}
\small

\[
    y_k
    =
    \mathbf{h}_k^H\mathbf{D}_k\mathbf{v}_k s_k
    +
    \sum_{j\ne k}^{K}
    \mathbf{h}_k^H\mathbf{D}_j\mathbf{v}_j s_j
    +
    n_k
\]

\begin{itemize}
    \item \(y_k\): signal received by user \(k\).
    \item \(\mathbf{h}_k\in\mathbb{C}^{N}\): channel vector from all \(N\) antennas to user \(k\).
    \item \((\cdot)^H\): Hermitian transpose.
    \item \(\mathbf{h}_k^H\mathbf{D}_k\mathbf{v}_k s_k\): desired signal for user \(k\).
    \item \(\sum_{j\ne k}^{K}\mathbf{h}_k^H\mathbf{D}_j\mathbf{v}_j s_j\): interference from other users.
    \item \(j\): index of another user, where \(j\ne k\).
    \item \(n_k\): additive noise at user \(k\).
\end{itemize}

\end{frame}

%------------------------------------------------
\begin{frame}{Layer 2: SINR of Each User}
\scriptsize

\[
\rho_k(\mathbf{V})
=
\frac{
\left|\mathbf{h}_k^H\mathbf{D}_k\mathbf{v}_k\right|^2
}{
\sum_{\substack{j=1\\j\ne k}}^{K}
\left|\mathbf{h}_k^H\mathbf{D}_j\mathbf{v}_j\right|^2
+
\sigma_k^2
}
\]

\begin{itemize}
    \item \(\rho_k(\mathbf{V})\): SINR of user \(k\).
    \item \(\mathbf{V}=[\mathbf{v}_1,\ldots,\mathbf{v}_K]\in\mathbb{C}^{N\times K}\): complete precoding matrix.
    \item Numerator:
    \[
        \left|\mathbf{h}_k^H\mathbf{D}_k\mathbf{v}_k\right|^2
    \]
    is the useful received signal power for user \(k\).

    \item Denominator:
    \[
        \sum_{j\ne k}
        \left|\mathbf{h}_k^H\mathbf{D}_j\mathbf{v}_j\right|^2
        +
        \sigma_k^2
    \]
    is interference power plus noise power.

    \item \(\sigma_k^2\): noise power at user \(k\). The paper assumes \(\sigma_k^2=1\) for simplicity.
\end{itemize}

\end{frame}

%------------------------------------------------
\begin{frame}{Layer 3: Finite Blocklength Rate Model}
\scriptsize

\[
R_k(\rho_k)
=
\log(1+\rho_k)
-
\frac{Q^{-1}(\epsilon_k)}{\sqrt{M}}
\sqrt{B(\rho_k)}
\]

\[
B(\rho_k)=1-(1+\rho_k)^{-2}
\]

\begin{itemize}
    \item \(R_k(\rho_k)\): achievable finite-blocklength rate of user \(k\).
    \item \(\rho_k\): SINR of user \(k\).
    \item \(\log(1+\rho_k)\): classical infinite-blocklength Shannon rate term.
    \item \(M\): blocklength, i.e., number of channel uses used to transmit a codeword.
    \item \(\epsilon_k\): target block error rate of user \(k\).
    \item \(Q^{-1}(\epsilon_k)\): inverse Gaussian \(Q\)-function; represents reliability penalty.
    \item \(B(\rho_k)\): channel dispersion term.
    \item The second term reduces the rate because short-packet transmission cannot achieve Shannon capacity.
\end{itemize}

\end{frame}

%------------------------------------------------
\begin{frame}{Layer 4: Mixed Power Constraints}
\small

\[
    \sum_{k=1}^{K}
    \mathbf{v}_k^H\mathbf{S}_i\mathbf{v}_k
    \le P_i,
    \quad \forall i=1,\ldots,L
\]

\begin{itemize}
    \item \(i\): index of a power constraint.
    \item \(L\): total number of power constraints.
    \item \(\mathbf{v}_k\): precoder for user \(k\).
    \item \(\mathbf{S}_i\in\mathbb{C}^{N\times N}\): positive semi-definite matrix selecting which antenna, group, or total power is constrained.
    \item \(P_i\): power budget of constraint \(i\).
    \item \(\mathbf{v}_k^H\mathbf{S}_i\mathbf{v}_k\): contribution of user \(k\)'s precoder to power constraint \(i\).
\end{itemize}

\end{frame}
%------------------------------------------------
\begin{frame}{Layer 4: Types of Power Constraints}
\small

\begin{itemize}
    \item The mixed power constraint framework can represent different practical power limits.

    \item \textbf{Sum Power Constraint (SPC):}
    \[
        \mathbf{S}_i=\mathbf{I}_N
    \]
    This limits the total transmit power across all antennas.

    \item \textbf{Per-Antenna Power Constraint (PAPC):}
    \[
        \mathbf{S}_i=\mathbf{e}_{m_i}\mathbf{e}_{m_i}^H
    \]
    This limits the transmit power of antenna \(m_i\).

    \item \textbf{Per-Group Power Constraint (PGPC):}
    \[
        \mathbf{S}_i=\sum_{m\in\mathcal{G}_i}
        \mathbf{e}_m\mathbf{e}_m^H
    \]
    This limits the total power of antenna group \(\mathcal{G}_i\).

\item \(m\): antenna index, where \(m=1,\ldots,N\), \([\mathbf{v}_k]_{m_i}\): the precoder coefficient applied at antenna \(m_i\) for user \(k\).
, \(m_i\): the specific antenna index selected by constraint \(i\), \(\mathbf{e}_{m_i}\): unit vector that selects antenna \(m_i\)..
\(\mathcal{G}_i\): set of antenna indices belonging to group \(i\).



\end{itemize}

\end{frame}

%------------------------------------------------
\begin{frame}{Layer 4: How \(\mathbf{S}_i\) Creates Different Power Constraints}
\scriptsize

\[
    \sum_{k=1}^{K}\mathbf{v}_k^H\mathbf{S}_i\mathbf{v}_k \le P_i
\]

\begin{itemize}
    \item \(\mathbf{S}_i\) acts like a power selection matrix.

    \item If \(\mathbf{S}_i=\mathbf{I}_N\):
    \[
        \sum_{k=1}^{K}\mathbf{v}_k^H\mathbf{v}_k \le P_i
    \]
    This becomes a total sum power constraint.

    \item If \(\mathbf{S}_i=\mathbf{e}_{m_i}\mathbf{e}_{m_i}^H\):
    \[
        \sum_{k=1}^{K}|[\mathbf{v}_k]_{m_i}|^2 \le P_i
    \]
    This limits only antenna \(m_i\).

    \item If \(\mathbf{S}_i=\sum_{m\in\mathcal{G}_i}\mathbf{e}_m\mathbf{e}_m^H\):
    \[
        \sum_{k=1}^{K}\sum_{m\in\mathcal{G}_i}|[\mathbf{v}_k]_m|^2 \le P_i
    \]
    This limits the total power of antenna group \(\mathcal{G}_i\).

    \item Thus, changing \(\mathbf{S}_i\) changes what part of the precoder power is measured.
\end{itemize}

\end{frame}
%------------------------------------------------
%------------------------------------------------
\begin{frame}{Layer 4: Normalized Mixed Power Constraint}
\scriptsize

\[
    \mathbf{X}_i=\frac{\mathbf{S}_i}{P_i}
\]

\[
    \sum_{k=1}^{K}\mathbf{v}_k^H\mathbf{X}_i\mathbf{v}_k \le 1,
    \quad \forall i=1,\ldots,L
\]

Since we want this power constraint for every constraint $i \in \{1,...,L\}$ 
\[
    \max_{1\le i\le L}
    \left\{
    \sum_{k=1}^{K}\mathbf{v}_k^H\mathbf{X}_i\mathbf{v}_k
    \right\}
    \le 1
\]

\begin{itemize}
    \item \(\mathbf{X}_i\): normalized power selection matrix for constraint \(i\).
    \item \(\mathbf{S}_i\): original power selection matrix.
    \item \(P_i\): power budget of constraint \(i\).
    \item Normalization converts every power constraint into the same upper bound \(1\).
    \item The maximum form means all power constraints must be satisfied simultaneously.
\end{itemize}

\end{frame}

%------------------------------------------------
%------------------------------------------------
\begin{frame}{Layer 4: Why the Maximum Form is Used}
\scriptsize

\[
    \sum_{k=1}^{K}\mathbf{v}_k^H\mathbf{S}_i\mathbf{v}_k \le P_i,
    \quad \forall i=1,\ldots,L
\]

\[
    \mathbf{X}_i=\frac{\mathbf{S}_i}{P_i}
\]

\[
    \sum_{k=1}^{K}\mathbf{v}_k^H\mathbf{X}_i\mathbf{v}_k \le 1,
    \quad \forall i=1,\ldots,L
\]

\[
    \max_{1\le i\le L}
    \left\{
    \sum_{k=1}^{K}\mathbf{v}_k^H\mathbf{X}_i\mathbf{v}_k
    \right\}
    \le 1
\]

\begin{itemize}
    \item The maximum form is used in Section II-B of the paper after defining the normalized matrix \(\mathbf{X}_i\).

    \item It converts \(L\) separate constraints into one compact expression.

    \item If the largest normalized power value is less than or equal to \(1\), then all individual power constraints are satisfied.

    \item This is useful because the original optimization problem can now write the mixed power constraints as one single constraint.

    \item \(\mathbf{S}_i\): original matrix selecting the antenna, antenna group, or total transmit power being constrained.
    \item \(L\): total number of power constraints.

\end{itemize}

\end{frame}

%------------------------------------------------
\begin{frame}{Layer 4: QoS Constraint}
\scriptsize

\[
    R_k(\rho_k(\mathbf{V})) \ge \widehat{R}_k,
    \quad \forall k=1,\ldots,K
\]

\[
    \rho_k(\mathbf{V}) \ge \widehat{\rho}_k,
    \quad \forall k=1,\ldots,K
\]

\[
    R_k(\widehat{\rho}_k)=\widehat{R}_k
\]

\begin{itemize}
    \item QoS means each user must receive at least a minimum required rate.
    \item \(R_k(\rho_k(\mathbf{V}))\): finite-blocklength rate achieved by user \(k\).
    \item \(\widehat{R}_k\): minimum required rate for user \(k\).
    \item \(\rho_k(\mathbf{V})\): SINR achieved by user \(k\) using precoding matrix \(\mathbf{V}\).
    \item \(\widehat{\rho}_k\): minimum SINR needed to achieve rate \(\widehat{R}_k\).
    \item Since \(R_k(\rho_k)\) increases with \(\rho_k\), the rate constraint can be rewritten as an SINR constraint.
\end{itemize}

\end{frame}

%------------------------------------------------
\begin{frame}{Problems with this QoS Transformation}
\scriptsize

\[
    R_k(\rho_k)
    =
    \log(1+\rho_k)
    -
    \frac{Q^{-1}(\epsilon_k)}{\sqrt{M}}
    \sqrt{B(\rho_k)}
\]

\[
    B(\rho_k)=1-(1+\rho_k)^{-2}
\]

\begin{itemize}
    \item The paper rewrites the QoS constraint:
    \[
        R_k(\rho_k(\mathbf{V}))\ge \widehat{R}_k
    \]
    as:
    \[
        \rho_k(\mathbf{V})\ge \widehat{\rho}_k
    \]

    \item This requires \(R_k(\rho_k)\) to be increasing with respect to \(\rho_k\).

    \item However, the dispersion term also increases with SINR:
    \[
        \frac{dB(\rho_k)}{d\rho_k}
        =
        2(1+\rho_k)^{-3}>0
    \]

    \item Therefore, increasing SINR increases both:
    \begin{itemize}
        \item the useful term \(\log(1+\rho_k)\),
        \item the penalty term \(\sqrt{B(\rho_k)}\).
    \end{itemize}

    \item The transformation is valid only when the useful rate increase dominates the penalty increase, ie: the shannon rate term increase dominates over  channel dispersion term.
    on the valid increasing region of the FBL rate curve.
\end{itemize}

\end{frame}
%------------------------------------------------
\begin{frame}{Layer 5: Original Optimization Problem}
\scriptsize

\[
\max_{\mathbf{V}}
\quad
\sum_{k=1}^{K} w_k R_k(\rho_k(\mathbf{V}))
\]

\[
\text{subject to}
\quad
\max_{1\le i\le L}
\left\{
\sum_{k=1}^{K}\mathbf{v}_k^H\mathbf{X}_i\mathbf{v}_k
\right\}
\le 1
\]

\[
\rho_k(\mathbf{V})\ge \widehat{\rho}_k,
\quad \forall k=1,\ldots,K
\]

\begin{itemize}
    \item \(\mathbf{V}=[\mathbf{v}_1,\ldots,\mathbf{v}_K]\): complete precoding matrix.
    \item \(w_k\): priority weight of user \(k\).
    \item \(R_k(\rho_k(\mathbf{V}))\): FBL rate of user \(k\).
    \item \(\mathbf{X}_i\): normalized power constraint matrix.
    \item \(\widehat{\rho}_k\): minimum required SINR for user \(k\).
    \item The problem is difficult because \(\rho_k(\mathbf{V})\) is non-convex in \(\mathbf{V}\).
\end{itemize}

\end{frame}

%------------------------------------------------
%------------------------------------------------
\begin{frame}{Layer 6: Why Directly Solving the Problem is Hard}
\scriptsize

\[
R_k(\rho_k)
=
\log(1+\rho_k)
-
\frac{Q^{-1}(\epsilon_k)}{\sqrt{M}}
\sqrt{1-(1+\rho_k)^{-2}}
\]

\[
\rho_k(\mathbf{V})
=
\frac{
\left|\mathbf{h}_k^H\mathbf{D}_k\mathbf{v}_k\right|^2
}{
\sum_{j\ne k}
\left|\mathbf{h}_k^H\mathbf{D}_j\mathbf{v}_j\right|^2
+
\sigma_k^2
}
\]

\begin{itemize}
    \item The difficulty comes from two sources.

    \item \textbf{First: SINR is non-convex in the precoders.}
    \[
        \rho_k(\mathbf{V})
    \]
    is a ratio of quadratic functions of the precoding vectors.

    \item Improving one user's signal power can increase interference to other users.

    \item \textbf{Second: the finite-blocklength rate is not the classical Shannon rate.}

    \item In the infinite-blocklength case, the rate is:
    \[
        \log(1+\rho_k)
    \]

    \item In the finite-blocklength case, the rate contains the extra dispersion penalty:
    \[
        \frac{Q^{-1}(\epsilon_k)}{\sqrt{M}}
        \sqrt{1-(1+\rho_k)^{-2}}
    \]

    \item Because of this penalty term, standard infinite-blocklength precoding tools such as closed-form WMMSE updates are not directly applicable.

    \item Therefore, the paper avoids directly optimizing \(\mathbf{V}\) and later reformulates the problem using target SINR allocation.
\end{itemize}

\end{frame}

%------------------------------------------------
\begin{frame}{Layer 6: Main Reformulation Idea}
\scriptsize

\[
    \gamma
    =
    [\gamma_1,\gamma_2,\ldots,\gamma_K]^T
\]

\begin{itemize}
    \item Instead of directly optimizing the precoders \(\mathbf{V}\), the paper optimizes target SINR values.

    \item \(\gamma_k\): target SINR assigned to user \(k\).

    \item \(\gamma\in\mathbb{R}^{K}\): vector containing the target SINRs of all users.

    \item The original problem asks:
    \[
        \text{What precoders } \mathbf{V} \text{ maximize the FBL weighted sum rate?}
    \]

    \item The reformulated problem asks:
    \[
        \text{What SINR targets } \gamma_k \text{ should users receive?}
    \]

    \item After choosing \(\gamma\), the paper checks whether there exists a precoding matrix \(\mathbf{V}\) that can achieve those SINR targets under the mixed power constraints.

    \item This separates the problem into:
    \[
        \text{SINR allocation}
        \quad+\quad
        \text{precoder feasibility checking}.
    \]
\end{itemize}

\end{frame}
%------------------------------------------------
\begin{frame}{Layer 6: Main Reformulation Idea}
\scriptsize

\[
    \gamma
    =
    [\gamma_1,\gamma_2,\ldots,\gamma_K]^T
\]

\begin{itemize}
    \item Instead of directly optimizing the precoders \(\mathbf{V}\), the paper optimizes target SINR values.

    \item \(\gamma_k\): target SINR assigned to user \(k\).

    \item \(\gamma\in\mathbb{R}^{K}\): vector containing the target SINRs of all users.

    \item The original problem asks:
    \[
        \text{What precoders } \mathbf{V} \text{ maximize the FBL weighted sum rate?}
    \]

    \item The reformulated problem asks:
    \[
        \text{What SINR targets } \gamma_k \text{ should users receive?}
    \]

    \item After choosing \(\gamma\), the paper checks whether there exists a precoding matrix \(\mathbf{V}\) that can achieve those SINR targets under the mixed power constraints.

    \item This separates the problem into:
    \[
        \text{SINR allocation}
        \quad+\quad
        \text{precoder feasibility checking}.
    \]
\end{itemize}

\end{frame}

\begin{frame}{Layer 6: Theorem 1 -- Equivalent SINR Allocation}
\scriptsize

\begin{block}{Theorem 1}
The original precoder optimization problem over \(\mathbf{V}\) is equivalent to an optimization problem over target SINRs \(\boldsymbol{\gamma}\).
\end{block}

\[
\max_{\mathbf{V}}
\sum_{k=1}^{K}w_kR_k(\rho_k(\mathbf{V}))
\quad
\Longleftrightarrow
\quad
\max_{\boldsymbol{\gamma}}
\sum_{k=1}^{K}w_kR_k(\gamma_k)
\]

\[
\text{subject to}
\quad
P(\boldsymbol{\gamma})\le 1,
\quad
\boldsymbol{\gamma}\ge \widehat{\boldsymbol{\rho}}
\]

\begin{itemize}
    \item \(\mathbf{V}\): original precoding matrix.
    \item \(\rho_k(\mathbf{V})\): actual SINR achieved by user \(k\).
    \item \(\boldsymbol{\gamma}\): target SINR vector chosen by the new problem.
    \item \(P(\boldsymbol{\gamma})\): minimum normalized power needed to achieve \(\boldsymbol{\gamma}\).
    \item \(P(\boldsymbol{\gamma})\le 1\): the target SINRs are feasible under the mixed power constraints.
    \item Theorem 1 is the key step that changes the problem from designing precoders directly to allocating feasible SINR targets.
\end{itemize}

\end{frame}
%------------------------------------------------
\begin{frame}{Layer 6: Definition of \(P(\boldsymbol{\gamma})\)}
\scriptsize

\[
P(\boldsymbol{\gamma})
=
\min_{\mathbf{V}}
\max_{1\le i\le L}
\left\{
\sum_{k=1}^{K}
\mathbf{v}_k^H\mathbf{X}_i\mathbf{v}_k
\right\}
\]

\[
\text{subject to}
\quad
\rho_k(\mathbf{V})\ge \gamma_k,
\quad \forall k=1,\ldots,K
\]

\begin{itemize}
    \item \(P(\boldsymbol{\gamma})\): minimum normalized power required to achieve the target SINR vector \(\boldsymbol{\gamma}\).
    \item \(\boldsymbol{\gamma}=[\gamma_1,\ldots,\gamma_K]^T\): target SINR vector.
    \item \(\gamma_k\): target SINR required for user \(k\).
    \item \(\mathbf{V}=[\mathbf{v}_1,\ldots,\mathbf{v}_K]\): precoding matrix being optimized.
    \item \(\mathbf{X}_i\): normalized matrix for power constraint \(i\).
    \item If \(P(\boldsymbol{\gamma})\le 1\), the SINR target is feasible.
    \item If \(P(\boldsymbol{\gamma})>1\), the SINR target requires more power than allowed.
\end{itemize}

\end{frame}

%------------------------------------------------
\begin{frame}{Layer 6: Why \(P(\boldsymbol{\gamma})\) is a Min--Max Problem}
\scriptsize

\[
P(\boldsymbol{\gamma})
=
\min_{\mathbf{V}}
\max_{1\le i\le L}
\left\{
\sum_{k=1}^{K}
\mathbf{v}_k^H\mathbf{X}_i\mathbf{v}_k
\right\}
\]

\[
\text{subject to}
\quad
\rho_k(\mathbf{V})\ge \gamma_k,
\quad \forall k=1,\ldots,K
\]

\begin{itemize}
    \item After choosing target SINRs
    \[
        \boldsymbol{\gamma}=[\gamma_1,\ldots,\gamma_K]^T,
    \]
    the paper must check whether some precoder \(\mathbf{V}\) can achieve them.

    \item The constraint
    \[
        \rho_k(\mathbf{V})\ge \gamma_k
    \]
    means user \(k\)'s achieved SINR must be at least its target SINR.

    \item The inner maximum appears because the system has \(L\) normalized mixed power constraints:
    \[
        \sum_{k=1}^{K}\mathbf{v}_k^H\mathbf{X}_i\mathbf{v}_k \le 1,
        \quad i=1,\ldots,L.
    \]

    \item Therefore,
    \[
        \max_{1\le i\le L}
        \left\{
        \sum_{k=1}^{K}\mathbf{v}_k^H\mathbf{X}_i\mathbf{v}_k
        \right\}
    \]
    measures the worst violated or most heavily used normalized power constraint.

    \item The outer minimization over \(\mathbf{V}\) searches for the best precoder that achieves \(\boldsymbol{\gamma}\) using the least possible normalized power.

    \item Hence, \(P(\boldsymbol{\gamma})\) means:
    \[
        \text{minimum worst-case normalized power needed to achieve }
        \boldsymbol{\gamma}.
    \]

    \item If
    \[
        P(\boldsymbol{\gamma})\le 1,
    \]
    then all mixed power constraints can be satisfied.

    \item If
    \[
        P(\boldsymbol{\gamma})>1,
    \]
    then the target SINRs require more power than the system allows.
\end{itemize}

\end{frame}

%------------------------------------------------
\begin{frame}{Layer 6: Why the Outer Minimization is Added}
\scriptsize

\[
P(\boldsymbol{\gamma})
=
\min_{\mathbf{V}}
\max_{1\le i\le L}
\left\{
\sum_{k=1}^{K}\mathbf{v}_k^H\mathbf{X}_i\mathbf{v}_k
\right\}
\]

\[
\text{subject to}
\quad
\rho_k(\mathbf{V})\ge \gamma_k,
\quad \forall k
\]

\begin{itemize}
    \item After choosing a target SINR vector \(\boldsymbol{\gamma}\), many different precoders \(\mathbf{V}\) may achieve it.

    \item The paper does not want to check only one arbitrary precoder.

    \item Instead, it asks:
    \[
        \text{Among all precoders that achieve } \boldsymbol{\gamma},
        \text{ which one uses the least worst-case normalized power?}
    \]

    \item This is why the outer minimization is added:
    \[
        \min_{\mathbf{V}}
    \]

    \item The inner maximum measures the worst power constraint for a chosen \(\mathbf{V}\).

    \item The outer minimum finds the best possible \(\mathbf{V}\) that minimizes that worst constraint.

    \item Therefore:
    \[
        P(\boldsymbol{\gamma})
        =
        \text{minimum worst-case normalized power needed to achieve }
        \boldsymbol{\gamma}.
    \]

    \item If \(P(\boldsymbol{\gamma})\le 1\), then at least one feasible precoder exists.
\end{itemize}

\end{frame}

%------------------------------------------------
\begin{frame}{Layer 7: Why Solving \(P(\boldsymbol{\gamma})\) is Still Difficult}
\scriptsize

\[
P(\boldsymbol{\gamma})
=
\min_{\mathbf{V}}
\max_{1\le i\le L}
\left\{
\sum_{k=1}^{K}
\mathbf{v}_k^H\mathbf{X}_i\mathbf{v}_k
\right\}
\]

\[
\text{subject to}
\quad
\rho_k(\mathbf{V})\ge\gamma_k,
\quad \forall k=1,\ldots,K
\]

\begin{itemize}
    \item Even after the SINR-allocation reformulation, computing \(P(\boldsymbol{\gamma})\) is not trivial.

    \item The objective contains a maximum over \(L\) mixed power constraints.

    \item The constraints still contain the SINR expression:
    \[
        \rho_k(\mathbf{V})
        =
        \frac{
        |\mathbf{h}_k^H\mathbf{D}_k\mathbf{v}_k|^2
        }{
        \sum_{j\ne k}
        |\mathbf{h}_k^H\mathbf{D}_j\mathbf{v}_j|^2+\sigma_k^2
        }.
    \]

    \item This SINR constraint couples all users because each \(\mathbf{v}_j\) can create interference for user \(k\).

    \item Therefore, the paper next transforms this power feasibility problem into an equivalent dual-domain convex problem.
\end{itemize}

\end{frame}

%------------------------------------------------
%------------------------------------------------
\begin{frame}{Layer 7: Converting \(P(\boldsymbol{\gamma})\) into a Dual Problem}
\scriptsize

\[
P(\boldsymbol{\gamma})
=
\min_{\mathbf{V}}
\max_{1\le i\le L}
\left\{
\sum_{k=1}^{K}
\mathbf{v}_k^H\mathbf{X}_i\mathbf{v}_k
\right\}
\]

\[
\text{subject to}
\quad
\rho_k(\mathbf{V})\ge\gamma_k,
\quad \forall k=1,\ldots,K
\]

\begin{itemize}
    \item The paper does not solve this constrained minimization directly.

    \item Instead, it converts Problem (11) into an equivalent dual-domain problem.

    \item New variables are introduced:
    \[
        \mathbf{q}=[q_1,\ldots,q_L]^T,
        \qquad
        \boldsymbol{\lambda}=[\lambda_1,\ldots,\lambda_K]^T
    \]

    \item \(q_i\): dual variable associated with power constraint \(i\).

    \item \(\lambda_k\): dual variable associated with the SINR constraint
    \[
        \rho_k(\mathbf{V})\ge\gamma_k.
    \]

    \item The goal is to compute \(P(\boldsymbol{\gamma})\) efficiently while accounting for both mixed power constraints and SINR target constraints.
\end{itemize}

\end{frame}

%------------------------------------------------
%------------------------------------------------
\begin{frame}{Layer 7: Lemma 1 -- Dual Form of \(P(\boldsymbol{\gamma})\)}
\scriptsize

\[
\max_{\mathbf{q},\boldsymbol{\lambda}}
\quad
\sum_{k=1}^{K}\lambda_k
\]

\[
\text{subject to}
\quad
\mathbf{A}_k(\mathbf{q},\boldsymbol{\lambda})\succeq 0,
\quad \forall k=1,\ldots,K
\]

\[
\sum_{i=1}^{L}q_i=1,
\quad
\mathbf{q}\ge 0,
\quad
\boldsymbol{\lambda}\ge 0
\]

\begin{itemize}
    \item Lemma 1 states that the difficult power problem \(P(\boldsymbol{\gamma})\) has the same optimal value as this dual problem.
    \item \(\mathbf{q}=[q_1,\ldots,q_L]^T\): dual variables for the \(L\) mixed power constraints.
    \item \(q_i\): weight assigned to power constraint \(i\).
    \item \(\boldsymbol{\lambda}=[\lambda_1,\ldots,\lambda_K]^T\): dual variables for the \(K\) SINR constraints.
    \item \(\lambda_k\): dual variable associated with user \(k\)'s SINR target \(\gamma_k\).
    \item \(\mathbf{A}_k(\mathbf{q},\boldsymbol{\lambda})\succeq0\): matrix feasibility condition for user \(k\).
    \item \(\succeq0\): positive semidefinite matrix condition.
\end{itemize}

\end{frame}
%------------------------------------------------
\begin{frame}{Layer 7: Meaning of \(\mathbf{A}_k(\mathbf{q},\boldsymbol{\lambda})\)}
\scriptsize

\[
\mathbf{A}_k(\mathbf{q},\boldsymbol{\lambda})
=
\sum_{i=1}^{L}q_i\mathbf{X}_i
+
\mathbf{D}_k
\left(
\sum_{j\ne k}\lambda_j\mathbf{h}_j\mathbf{h}_j^H
-
\frac{\lambda_k}{\gamma_k}
\mathbf{h}_k\mathbf{h}_k^H
\right)
\mathbf{D}_k
\]

\[
\mathbf{A}_k(\mathbf{q},\boldsymbol{\lambda})\succeq0,
\quad \forall k
\]

\begin{itemize}
    \item \(\mathbf{A}_k(\mathbf{q},\boldsymbol{\lambda})\): matrix condition for user \(k\) in the dual problem.
    \item \(\sum_{i=1}^{L}q_i\mathbf{X}_i\): weighted combination of all normalized power constraints.
    \item \(\lambda_j\mathbf{h}_j\mathbf{h}_j^H\): interference-related term from user \(j\).
    \item \(\frac{\lambda_k}{\gamma_k}\mathbf{h}_k\mathbf{h}_k^H\): desired-signal term for user \(k\), scaled by its SINR target.
    \item \(\mathbf{D}_k\): selects antennas that serve user \(k\).
    \item \(\mathbf{A}_k\succeq0\) ensures the dual function remains bounded.
\end{itemize}

\end{frame}
%------------------------------------------------
\begin{frame}{Layer 7: Where \(\mathbf{A}_k\) Comes From}
\scriptsize

\[
p_i(\mathbf{V})
=
\sum_{k=1}^{K}
\mathbf{v}_k^H\mathbf{X}_i\mathbf{v}_k
\]

\[
\max_i p_i(\mathbf{V})
\quad \Longrightarrow \quad
p_i(\mathbf{V})\le \xi,\quad \forall i
\]

\[
\mathcal{L}(\xi,\mathbf{V},\mathbf{q},\boldsymbol{\lambda})
=
\xi
+
\sum_{i=1}^{L}q_i(p_i(\mathbf{V})-\xi)
-
\sum_{k=1}^{K}\lambda_k
\left(
\frac{1}{\gamma_k}
|\mathbf{h}_k^H\mathbf{D}_k\mathbf{v}_k|^2
-
\sum_{j\ne k}
|\mathbf{h}_k^H\mathbf{D}_j\mathbf{v}_j|^2
-
1
\right)
\]

\begin{itemize}
    \item \(\xi\): slack variable replacing the maximum power term.
    \item \(p_i(\mathbf{V})\): normalized power used under constraint \(i\).
    \item \(q_i\): dual variable for \(p_i(\mathbf{V})\le\xi\).
    \item \(\lambda_k\): dual variable for user \(k\)'s SINR constraint.
    \item After collecting all terms containing each \(\mathbf{v}_k\), the quadratic coefficient becomes \(\mathbf{A}_k(\mathbf{q},\boldsymbol{\lambda})\).
\end{itemize}

\end{frame}
%------------------------------------------------

%------------------------------------------------
\begin{frame}{Layer 7: From Lagrangian to the Dual Problem}
\scriptsize

\[
\mathcal{L}
=
\sum_{k=1}^{K}\lambda_k
-
\xi\left(\sum_{i=1}^{L}q_i-1\right)
+
\sum_{k=1}^{K}
\mathbf{v}_k^H
\mathbf{A}_k(\mathbf{q},\boldsymbol{\lambda})
\mathbf{v}_k
\]

\[
\sum_{i=1}^{L}q_i=1,
\qquad
\mathbf{A}_k(\mathbf{q},\boldsymbol{\lambda})\succeq0,
\quad \forall k
\]

\[
\Longrightarrow
\quad
\max_{\mathbf{q},\boldsymbol{\lambda}}
\sum_{k=1}^{K}\lambda_k
\]

\[
\text{subject to}
\quad
\mathbf{A}_k(\mathbf{q},\boldsymbol{\lambda})\succeq0,\ \forall k,
\quad
\sum_{i=1}^{L}q_i=1,
\quad
\mathbf{q}\ge0,
\quad
\boldsymbol{\lambda}\ge0
\]

\begin{itemize}
    \item The dual problem keeps only values of \(\mathbf{q}\) and \(\boldsymbol{\lambda}\) that make the Lagrangian bounded.
    \item \(\sum_i q_i=1\): removes the unbounded slack-variable term involving \(\xi\).
    \item \(\mathbf{A}_k\succeq0\): prevents the quadratic precoder term from going to \(-\infty\).
    \item After these conditions hold, minimizing over \(\xi\) and \(\mathbf{V}\) leaves the dual objective \(\sum_k\lambda_k\).
\end{itemize}

\end{frame}

%------------------------------------------------
\begin{frame}{Layer 8: Solving the Dual Problem -- Updating \(\boldsymbol{\lambda}\)}
\scriptsize

\[
\lambda_k
\leftarrow
\frac{\gamma_k}
{\mathbf{h}_k^H\mathbf{D}_k\mathbf{C}_k^{-1}\mathbf{D}_k\mathbf{h}_k},
\quad \forall k=1,\ldots,K
\]

\[
\mathbf{C}_k
=
\mathbf{D}_k
\left(
\sum_{j\ne k}
\lambda_j\mathbf{h}_j\mathbf{h}_j^H
\right)
\mathbf{D}_k
+
\sum_{i=1}^{L}q_i\mathbf{X}_i
\]

\begin{itemize}
    \item This update is used inside Algorithm 1 to solve the dual problem from Lemma 1.

    \item \(\lambda_k\): dual variable associated with user \(k\)'s SINR target constraint.

    \item \(\gamma_k\): target SINR assigned to user \(k\).

    \item \(\mathbf{C}_k\): effective interference-plus-power matrix seen while updating user \(k\)'s dual variable.

    \item \(\sum_{j\ne k}\lambda_j\mathbf{h}_j\mathbf{h}_j^H\): weighted interference contribution from all other users.

    \item \(\sum_{i=1}^{L}q_i\mathbf{X}_i\): weighted mixed-power-constraint contribution.

    \item The denominator measures how strongly user \(k\)'s channel can support its SINR target under the current interference and power weights.

    \item The update increases \(\lambda_k\) when achieving \(\gamma_k\) is harder.
\end{itemize}

\end{frame}
%------------------------------------------------
\begin{frame}{Layer 8: Reducing the Cost of Updating \(\boldsymbol{\lambda}\)}
\scriptsize

\[
\lambda_k
\leftarrow
\frac{\gamma_k}
{\mathbf{h}_k^H\mathbf{D}_k\mathbf{C}_k^{-1}\mathbf{D}_k\mathbf{h}_k}
\]

\[
\mathbf{C}_k
=
\mathbf{D}_k
\left(
\sum_{j\ne k}\lambda_j\mathbf{h}_j\mathbf{h}_j^H
\right)
\mathbf{D}_k
+
\sum_{i=1}^{L}q_i\mathbf{X}_i
\]

\begin{itemize}
    \item Directly computing \(\mathbf{C}_k^{-1}\) for every user is expensive.

    \item A full matrix inverse has complexity:
    \[
        \mathcal{O}(N^3)
    \]

    \item Since the update is needed for all \(K\) users, this can become computationally costly.

    \item The paper avoids repeated full inversions by updating the inverse recursively using the matrix inversion lemma.

    \item This reduces the complexity to approximately:
    \[
        \mathcal{O}(KN^2)
    \]

    \item \(N\): number of transmit antennas.
    \item \(K\): number of users.
    \item \(\mathbf{C}_k^{-1}\): inverse matrix needed to update \(\lambda_k\).
\end{itemize}

\end{frame}
%------------------------------------------------
\begin{frame}{Layer 8: Matrix Inversion Lemma Update}
\scriptsize

\[
\mathbf{C}_{k+1}
=
\mathbf{C}_k
-
\lambda_{k+1}\mathbf{D}_{k+1}\mathbf{h}_{k+1}\mathbf{h}_{k+1}^{H}\mathbf{D}_{k+1}
+
\lambda_k\mathbf{D}_k\mathbf{h}_k\mathbf{h}_k^{H}\mathbf{D}_k
\]

\[
\mathbf{C}_{k+1}^{-1}
=
\mathbf{C}_k^{-1}
-
\mathbf{C}_k^{-1}\mathbf{U}_k
\mathbf{G}_k^{-1}
\mathbf{W}_k
\mathbf{C}_k^{-1}
\]

\begin{itemize}
    \item The paper avoids recomputing \(\mathbf{C}_{k+1}^{-1}\) from scratch.

    \item \(\mathbf{C}_k\): matrix used when updating \(\lambda_k\).

    \item \(\mathbf{C}_{k+1}\): next matrix used when updating \(\lambda_{k+1}\).

    \item Since \(\mathbf{C}_{k+1}\) differs from \(\mathbf{C}_k\) by low-rank terms, its inverse can be updated efficiently.

    \item \(\mathbf{U}_k\), \(\mathbf{W}_k\), and \(\mathbf{G}_k\): auxiliary matrices used in the matrix inversion lemma.

    \item This reduces repeated inverse computation and makes Algorithm 1 computationally efficient.
\end{itemize}

\end{frame}
%------------------------------------------------
%------------------------------------------------
\begin{frame}{Layer 8: Updating \(\mathbf{q}\) with Projected Subgradient Ascent}
\scriptsize

\[
\widetilde{\mathbf{q}}^{(n)}
=
\mathbf{q}^{(n)}
+
\eta^{(n)}
\left(
\mathbf{q}^{(n)}-\mathbf{q}^{(n-1)}
\right)
\]

\[
\mathbf{q}^{(n+1)}
=
P_{\mathcal{Q}}
\left(
\widetilde{\mathbf{q}}^{(n)}
+
b^{(n)}
\mathbf{g}(\widetilde{\mathbf{q}}^{(n)})
\right)
\]

\[
\mathcal{Q}
=
\left\{
\mathbf{q}\mid
\sum_{i=1}^{L}q_i=1,\;
\mathbf{q}\ge0
\right\}
\]

\[
P_{\mathcal{Q}}(\mathbf{a})
=
\arg\min_{\mathbf{q}\in\mathcal{Q}}
\|\mathbf{q}-\mathbf{a}\|_2^2
\]

\begin{itemize}
    \item \(\mathbf{q}\): dual weight vector for the mixed power constraints.
    \item \(q_i\): weight assigned to power constraint \(i\).
    \item \(\widetilde{\mathbf{q}}^{(n)}\): momentum-based extrapolated value of \(\mathbf{q}\).
    \item \(\eta^{(n)}\): momentum coefficient that uses the previous movement direction.
    \item \(b^{(n)}\): subgradient step size.
    \item \(\mathbf{g}(\widetilde{\mathbf{q}}^{(n)})\): subgradient showing which power constraints are heavily used.
    \item \(\mathcal{Q}\): feasible set requiring \(q_i\ge0\) and \(\sum_i q_i=1\).
    \item \(P_{\mathcal{Q}}(\cdot)\): moves the updated vector back to the closest valid \(\mathbf{q}\).
    \item This is needed because the subgradient step may produce negative entries or make \(\sum_iq_i\ne1\).
\end{itemize}

\end{frame}
%------------------------------------------------

%------------------------------------------------
\begin{frame}{Layer 8: Subgradient Used to Update \(\mathbf{q}\)}
\scriptsize

\[
g_i(\widetilde{\mathbf{q}}^{(n)})
=
\sum_{k=1}^{K}
\left(\mathbf{v}_k^{(n)}\right)^H
\mathbf{X}_i
\mathbf{v}_k^{(n)}
\]

\[
\mathbf{v}_k
=
\sqrt{\delta_k}
\mathbf{C}_k^{-1}
\mathbf{D}_k\mathbf{h}_k
\]

\begin{itemize}
    \item \(g_i(\widetilde{\mathbf{q}}^{(n)})\): subgradient component for power constraint \(i\).

    \item It measures how much normalized power is used under constraint \(i\).

    \item \(\mathbf{v}_k^{(n)}\): precoder for user \(k\) at inner iteration \(n\).

    \item \(\mathbf{X}_i\): normalized matrix for power constraint \(i\).

    \item \(\delta_k\): power coefficient used to construct the precoder.

    \item \(\mathbf{C}_k^{-1}\mathbf{D}_k\mathbf{h}_k\): beamforming direction for user \(k\).

    \item The update increases \(q_i\) when constraint \(i\) is heavily used.
\end{itemize}

\end{frame}

%------------------------------------------------
\begin{frame}{Layer 8: Algorithm 1 -- Computing \(P(\boldsymbol{\gamma})\)}
\scriptsize

\begin{enumerate}
    \item Initialize \(\mathbf{q}^{(0)}\).

    \item Compute the extrapolated point:
    \[
        \widetilde{\mathbf{q}}^{(n)}
        =
        \mathbf{q}^{(n)}
        +
        \eta^{(n)}
        \left(\mathbf{q}^{(n)}-\mathbf{q}^{(n-1)}\right)
    \]

    \item Update \(\boldsymbol{\lambda}\) using the fixed-point rule:
    \[
        \lambda_k
        \leftarrow
        \frac{\gamma_k}
        {\mathbf{h}_k^H\mathbf{D}_k\mathbf{C}_k^{-1}\mathbf{D}_k\mathbf{h}_k}
    \]

    \item Construct the precoders:
    \[
        \mathbf{v}_k
        =
        \sqrt{\delta_k}
        \mathbf{C}_k^{-1}
        \mathbf{D}_k\mathbf{h}_k
    \]

    \item Update \(\mathbf{q}\) by projected subgradient ascent.
\end{enumerate}

\begin{itemize}
    \item Output: the value \(P(\boldsymbol{\gamma})\), the dual variables \(\mathbf{q}^{\star}\), \(\boldsymbol{\lambda}^{\star}\), and the corresponding precoders \(\mathbf{v}_k^{\star}\).
    \item This algorithm is used every time the outer algorithm needs to check whether a target SINR vector \(\boldsymbol{\gamma}\) is feasible.
\end{itemize}

\end{frame}


%------------------------------------------------
\begin{frame}{Layer 10: BDGM -- Main Idea}
\scriptsize

\[
\mathbf{z}
=
\log(1+\boldsymbol{\gamma})
\]

\[
z_k=\log(1+\gamma_k),
\quad k=1,\ldots,K
\]

\begin{itemize}
    \item After Theorem 1, the paper must solve the SINR-allocation problem:
    \[
        \max_{\boldsymbol{\gamma}}
        \sum_{k=1}^{K}w_kR_k(\gamma_k)
    \]
    subject to:
    \[
        P(\boldsymbol{\gamma})\le1,
        \quad
        \boldsymbol{\gamma}\ge\widehat{\boldsymbol{\rho}}.
    \]

    \item BDGM stands for \textbf{Bi-Direction Gradient Method}.

    \item Instead of updating \(\boldsymbol{\gamma}\) directly, BDGM updates:
    \[
        \mathbf{z}=\log(1+\boldsymbol{\gamma}).
    \]

    \item \(\mathbf{z}=[z_1,\ldots,z_K]^T\): rate-domain SINR variable.

    \item If the current point is feasible, BDGM moves forward to increase the weighted sum rate.

    \item If the current point is infeasible, BDGM moves backward to reduce the required power.
\end{itemize}

\end{frame}
%------------------------------------------------
\begin{frame}{Layer 10: BDGM Forward Step}
\scriptsize

\[
P_z(\mathbf{z}^{(t)})\le 1
\]

\[
\mathbf{z}^{(t+1)}
=
\mathbf{z}^{(t)}
+
\alpha^{(t)}
\nabla f_z(\mathbf{z}^{(t)})
\]

\[
[\nabla f_z(\mathbf{z})]_k
=
w_k
\left(
1
-
\frac{Q^{-1}(\epsilon_k)}{\sqrt{M}}
\frac{e^{-2z_k}}{\sqrt{1-e^{-2z_k}}}
\right)
\]

\begin{itemize}
    \item This step is used when the current SINR target is feasible.
    \item \(t\): outer iteration index.
    \item \(\mathbf{z}^{(t)}\): current rate-domain SINR vector.
    \item \(P_z(\mathbf{z})=P(e^{\mathbf{z}}-\mathbf{1})\): required normalized power in the \(z\)-domain.
    \item \(f_z(\mathbf{z})=f(e^{\mathbf{z}}-\mathbf{1})\): weighted sum-rate objective in the \(z\)-domain.
    \item \(\alpha^{(t)}\): forward step size.
    \item \(\nabla f_z(\mathbf{z}^{(t)})\): direction that increases weighted sum rate.
    \item Intuition: if the current point is feasible, BDGM increases the users' SINR targets to improve rate.
\end{itemize}

\end{frame}

%------------------------------------------------
\begin{frame}{Layer 10: BDGM Backward Step}
\scriptsize

\[
P_z(\mathbf{z}^{(t)})>1
\]

\[
\mathbf{z}^{(t+1)}
=
\max
\left(
\mathbf{z}^{(t)}
-
\beta^{(t)}
\nabla P_z(\mathbf{z}^{(t)}),
\log(1+\widehat{\boldsymbol{\rho}})
\right)
\]

\[
[\nabla P_z(\mathbf{z})]_k
=
e^{z_k}
\frac{\partial P(\boldsymbol{\gamma})}{\partial \gamma_k}
\Bigg|_{\boldsymbol{\gamma}=e^{\mathbf{z}}-\mathbf{1}}
\]

\begin{itemize}
    \item This step is used when the current SINR target is infeasible.
    \item \(P_z(\mathbf{z}^{(t)})>1\): current target SINRs require more normalized power than allowed.
    \item \(\beta^{(t)}\): backward step size.
    \item \(\nabla P_z(\mathbf{z}^{(t)})\): direction in which required power increases fastest.
    \item Moving in the negative gradient direction reduces the required power.
    \item \(\log(1+\widehat{\boldsymbol{\rho}})\): lower bound that prevents violating QoS.
    \item The max operation ensures the updated SINR target stays above the QoS threshold.
\end{itemize}

\end{frame}

%------------------------------------------------
\begin{frame}{Layer 10: BDGM Overall Procedure}
\scriptsize

\[
\mathbf{z}=\log(1+\boldsymbol{\gamma})
\]

\begin{enumerate}
    \item Start with an initial SINR target vector:
    \[
        \boldsymbol{\gamma}^{(0)}
    \]

    \item Convert it to the rate-domain variable:
    \[
        \mathbf{z}^{(0)}=\log(1+\boldsymbol{\gamma}^{(0)})
    \]

    \item At iteration \(t\), compute:
    \[
        P_z(\mathbf{z}^{(t)})
    \]

    \item If:
    \[
        P_z(\mathbf{z}^{(t)})\le1,
    \]
    use the forward step to increase the objective.

    \item If:
    \[
        P_z(\mathbf{z}^{(t)})>1,
    \]
    use the backward step to recover feasibility.
\end{enumerate}

\begin{itemize}
    \item BDGM alternates between rate improvement and power-feasibility correction.
    \item The output is a SINR target vector that satisfies QoS and mixed power constraints.
\end{itemize}

\end{frame}

%------------------------------------------------
\begin{frame}{Layer 10: BDGM Convergence}
\scriptsize

\[
P_z(\mathbf{z})\le1
\quad \text{means feasible}
\]

\[
P_z(\mathbf{z})>1
\quad \text{means infeasible}
\]

\[
\nabla f_z(\mathbf{z}^{\star})
=
\mu \nabla P_z(\mathbf{z}^{\star}),
\quad \mu\ge0
\]

\begin{itemize}
    \item The paper proves that BDGM converges to a point satisfying the KKT conditions of the SINR-allocation problem.

    \item \(\mathbf{z}^{\star}\): limiting point reached by BDGM.

    \item \(\nabla f_z(\mathbf{z}^{\star})\): gradient of the weighted sum-rate objective.

    \item \(\nabla P_z(\mathbf{z}^{\star})\): gradient of the required normalized power function.

    \item \(\mu\): non-negative multiplier for the feasibility boundary \(P_z(\mathbf{z})\le1\).

    \item At convergence, the forward rate-improvement direction and backward feasibility-correction direction balance each other.

    \item This means BDGM stops at a stationary feasible SINR allocation.
\end{itemize}

\end{frame}

%------------------------------------------------
\begin{frame}{Layer 11: Local Linearization of \(P(\boldsymbol{\gamma})\le1\)}
\scriptsize

\[
P(\boldsymbol{\gamma}^{(t)})
+
\nabla P(\boldsymbol{\gamma}^{(t)})^T
\left(
\boldsymbol{\gamma}-\boldsymbol{\gamma}^{(t)}
\right)
\le 1
\]

\begin{itemize}
    \item IWF approximates the difficult constraint \(P(\boldsymbol{\gamma})\le1\) around the current point \(\boldsymbol{\gamma}^{(t)}\).

    \item \(\boldsymbol{\gamma}^{(t)}\): SINR target vector at iteration \(t\).

    \item \(\boldsymbol{\gamma}\): candidate SINR target vector for the next update.

    \item \(P(\boldsymbol{\gamma}^{(t)})\): required normalized power at the current SINR target.

    \item \(\nabla P(\boldsymbol{\gamma}^{(t)})\): gradient of required power with respect to SINR targets.

    \item The linearized constraint estimates how the required power changes when moving from \(\boldsymbol{\gamma}^{(t)}\) to \(\boldsymbol{\gamma}\).

    \item This turns the difficult feasibility constraint into a simpler linear constraint for the next update.
\end{itemize}

\end{frame}
%------------------------------------------------
%------------------------------------------------
%------------------------------------------------
\begin{frame}{Layer 11: IWF Water-Filling Optimality Condition}
\scriptsize

\[
c_k^{(t)}
=
\left[\nabla P(\boldsymbol{\gamma}^{(t)})\right]_k
\]

\[
w_k
\frac{\partial R_k(\widehat{\gamma}_k^{(t+1)})}
{\partial \gamma_k}
=
\nu c_k^{(t)},
\quad \forall k
\]
\[
g^{(t)}
=
1
-
P(\boldsymbol{\gamma}^{(t)})
+
\nabla P(\boldsymbol{\gamma}^{(t)})^T
\boldsymbol{\gamma}^{(t)}
\]

\[
    \nabla P(\boldsymbol{\gamma}^{(t)})^T
    \widehat{\boldsymbol{\gamma}}^{(t+1)}
    =
    g^{(t)}
\]

\begin{itemize}
    \item \(g^{(t)}\): available linearized power budget at iteration \(t\).
    \item \(c_k^{(t)}\): marginal power cost of increasing user \(k\)'s SINR target at iteration \(t\).
    \item \(\nu\ge0\): dual variable, interpreted as the water level.
    \item \(w_k\): priority weight of user \(k\).
    \item \(\frac{\partial R_k}{\partial\gamma_k}\): marginal FBL rate gain from increasing \(\gamma_k\).
    \item The condition balances marginal rate gain against marginal power cost.
    \item \(\nu\) is found by bisection so that the linearized power constraint is satisfied.
\end{itemize}

\end{frame}
%------------------------------------------------
%------------------------------------------------
\begin{frame}{Layer 11: Why Gradients Balance at the IWF Optimum}
\scriptsize

\[
\max_{\boldsymbol{\gamma}\ge\widehat{\boldsymbol{\rho}}}
f(\boldsymbol{\gamma})
\quad
\text{subject to}
\quad
\sum_{k=1}^{K}c_k^{(t)}\gamma_k \le g^{(t)}
\]

\[
c_k^{(t)}
=
[\nabla P(\boldsymbol{\gamma}^{(t)})]_k
\]

\[
\mathcal{L}
=
f(\boldsymbol{\gamma})
-
\nu
\left(
\sum_{k=1}^{K}c_k^{(t)}\gamma_k-g^{(t)}
\right)
\]

\[
\frac{\partial \mathcal{L}}{\partial \gamma_k}=0
\]

\[
w_k\frac{\partial R_k(\gamma_k)}{\partial \gamma_k}
=
\nu c_k^{(t)}
\]

\begin{itemize}
    \item \(f(\boldsymbol{\gamma})=\sum_{k=1}^{K}w_kR_k(\gamma_k)\): weighted FBL sum-rate.
    \item The equality comes from setting the derivative of the Lagrangian to zero.
    \item It means marginal weighted rate gain equals marginal power cost multiplied by the common price \(\nu\).
\end{itemize}

\end{frame}
%------------------------------------------------
%------------------------------------------------
\begin{frame}{Layer 11: Smoothed IWF Update}
\scriptsize

\[
\boldsymbol{\gamma}^{(t+1)}
=
\eta^{(t)}
\widehat{\boldsymbol{\gamma}}^{(t+1)}
+
\left(1-\eta^{(t)}\right)
\boldsymbol{\gamma}^{(t)}
\]

\begin{itemize}
    \item \(\widehat{\boldsymbol{\gamma}}^{(t+1)}\): candidate SINR vector obtained from the water-filling subproblem.

    \item \(\boldsymbol{\gamma}^{(t)}\): current SINR target vector.

    \item \(\boldsymbol{\gamma}^{(t+1)}\): actual next SINR target vector used by the algorithm.

    \item \(\eta^{(t)}\in(0,1]\): smoothing parameter or step size.

    \item If \(\eta^{(t)}=1\), the algorithm fully moves to the water-filling solution.

    \item If \(\eta^{(t)}<1\), the algorithm moves only partially toward the candidate solution.

    \item This smoothing improves stability and helps maintain feasibility.
\end{itemize}

\end{frame}
%------------------------------------------------
\begin{frame}{Layer 11: IWF Overall Algorithm}
\scriptsize

\begin{enumerate}
    \item Initialize the SINR target vector:
    \[
        \boldsymbol{\gamma}^{(0)}
    \]

    \item At iteration \(t\), compute:
    \[
        P(\boldsymbol{\gamma}^{(t)})
        \quad \text{and} \quad
        \nabla P(\boldsymbol{\gamma}^{(t)})
    \]

    \item Locally linearize the power feasibility constraint:
    \[
        P(\boldsymbol{\gamma})\le1
    \]

    \item Solve the water-filling subproblem to obtain:
    \[
        \widehat{\boldsymbol{\gamma}}^{(t+1)}
    \]

    \item Apply the smoothed update:
    \[
        \boldsymbol{\gamma}^{(t+1)}
        =
        \eta^{(t)}
        \widehat{\boldsymbol{\gamma}}^{(t+1)}
        +
        (1-\eta^{(t)})
        \boldsymbol{\gamma}^{(t)}
    \]
\end{enumerate}

\begin{itemize}
    \item \(P(\boldsymbol{\gamma}^{(t)})\): required normalized power at the current SINR target.
    \item \(\nabla P(\boldsymbol{\gamma}^{(t)})\): marginal power cost of increasing SINR targets.
    \item \(\widehat{\boldsymbol{\gamma}}^{(t+1)}\): candidate water-filling SINR vector.
    \item \(\eta^{(t)}\): smoothing parameter.
\end{itemize}

\end{frame}


%------------------------------------------------
%------------------------------------------------
%------------------------------------------------
%------------------------------------------------
%------------------------------------------------
%------------------------------------------------
%------------------------------------------------
%------------------------------------------------
%------------------------------------------------
%------------------------------------------------
%------------------------------------------------
%------------------------------------------------
%------------------------------------------------
%------------------------------------------------
%------------------------------------------------
%------------------------------------------------
%------------------------------------------------
%------------------------------------------------
%------------------------------------------------
%------------------------------------------------
%------------------------------------------------
%------------------------------------------------
%------------------------------------------------
%------------------------------------------------
%------------------------------------------------
%------------------------------------------------

%------------------------------------------------
%------------------------------------------------

%------------------------------------------------
%------------------------------------------------
%------------------------------------------------
%------------------------------------------------
%------------------------------------------------
%------------------------------------------------
%------------------------------------------------
%------------------------------------------------
%------------------------------------------------
%------------------------------------------------
%------------------------------------------------
%------------------------------------------------
%------------------------------------------------
%------------------------------------------------
%------------------------------------------------
%------------------------------------------------
%------------------------------------------------
%------------------------------------------------

%------------------------------------------------
%------------------------------------------------
%------------------------------------------------
%------------------------------------------------
%------------------------------------------------
%------------------------------------------------
\end{document}